\section{Related Work}
\label{sec:related}
Table~\ref{tab:related-comparison} summarizes the selected systems discussed below. Their different tasks, interfaces, and evaluation settings allow a qualitative comparison, not a ranking or head-to-head benchmark.

\begin{table*}[!t]
\centering
\caption{Representative related work and the scope of its reported evaluation. The studies address different tasks and are not a common benchmark. OPC UA: Open Platform Communications Unified Architecture; PLC: programmable logic controller.}
\label{tab:related-comparison}
\footnotesize
\setlength{\tabcolsep}{5pt}
\begin{tabularx}{\textwidth}{@{}>{\raggedright\arraybackslash}p{0.16\textwidth}>{\raggedright\arraybackslash}p{0.40\textwidth}X@{}}
\toprule
Study & Main focus & Reported setting \\
\midrule
Biskupovic et al.~\cite{biskupovic2025supervision} & Agent-based industrial service selection and reselection & Service adaptability and performance in an industrial cyber--physical system \\
K\"{o}hle et al.~\cite{koehle2026llmdriven} & Module-specific agents and skill-based manufacturing & Nine simulated production challenges \\
Graul et al.~\cite{graul2026aas} & Grounded proposals, documented authorization, and proposed PLC-side checks & Two proof-of-concept scenarios using a simulated OPC UA server \\
Schall~\cite{schall2026safe} & Process-diagram, supervisory control and data acquisition (SCADA), and procedure-grounded verification & Five simulated scenarios and proposal evaluation \\
Vyas et al.~\cite{vyas2026detection} & Multi-agent fault-recovery planning with a process-plant twin & Simulated mixing and stirred-tank reactor studies \\
This work & Scene-bound signal mappings, shared dispatch, and route-specific governance & Simulator and mock-driver probes, scripted policy runs, and exploratory agent-driven episodes \\
\bottomrule
\end{tabularx}
\end{table*}

\subsection{Industrial Agents and Plant Integration}
Biskupovic et al. study online service selection and reselection in industrial cyber--physical systems~\cite{biskupovic2025supervision}. Xia et al. present a digital-twin-connected automation architecture and control generation~\cite{xia2025architecture,xia2025control}; other work generates programmable logic controller (PLC) programs~\cite{liu2026agents4plc}. In related control research, human instructions have also been incorporated into model-predictive control~\cite{wu2025instructmpc}.

Li et al. use structured instructions, retrieval, and iterative verification to generate and check finite-element scripts for Abaqus; that study evaluates engineering-artifact generation rather than runtime authority over process writes~\cite{li2026modsolagent}. Other studies examine large language model (LLM) agents, manufacturing skills, and natural-language machine interaction~\cite{koehle2026llmdriven,lim2025llmenabled,zhao2026shopfloor,hofmann2025nlc}. Vyas et al. combine multi-agent fault-recovery planning and plant-knowledge retrieval with Digital Process Plant Twin models of mixing and continuous stirred-tank reactors (CSTRs)~\cite{vyas2026detection}.

These studies address agent organization, generated actions, or task execution; the present study introduces no agent policy or control law but examines the runtime boundary that carries production context and direction-specific signal mappings into an operation and applies the required checks before dispatch.

\subsection{Industrial Interfaces and Semantic Commissioning}
Agent-facing interface work includes Model Context Protocol (MCP) access to manufacturing resources~\cite{silva2025mcpmanufacturing}, semantic recipes built from Open Platform Communications Unified Architecture (OPC UA) skills~\cite{koehle2025semanticrecipes}, industrial-protocol adapters~\cite{xavier2026industriConnect}, and tools for binding manufacturing data to a cognitive digital twin~\cite{liu2026semanticcommissioning}. These systems address capability discovery, protocol access, or commissioning. Here, configuration binds known protocol drivers to scene-specific read and write nodes, and a named operation resolves those bindings through a runtime path that checks authority and write limits; neither MCP nor a semantic node name is an authorization mechanism.

\subsection{Governed Operations and Verification}
Prior work also examines plant-facing verification. Graul et al. propose grounding plans in the Asset Administration Shell, with documented human authorization and PLC-side validation of OPC UA Methods~\cite{graul2026aas}. Schall verifies proposals against piping and instrumentation diagrams, SCADA data, and standard operating procedures; Lin and Chen describe a policy-governed sensor--manufacturing execution system (MES)--enterprise resource planning (ERP) digital thread~\cite{schall2026safe,lin2026agentic}. PLCBench evaluates autonomous agents against commercial PLCs with hardware-in-the-loop process simulation~\cite{zhou2026plcbench}. These studies echo long-standing authorization and safety concerns. Our implementation asks how configured signal bindings, route-specific checks, shared dispatch, and separate command and process records operate together at the write boundary. The contribution is system integration; the studies use different settings, and none establishes comparative safety.

\subsection{Production Context and Outcome Evidence}
MCP standardizes tool discovery and invocation~\cite{mcp2026}; it does not assign plant write authority. ISA-88 and ISA-95 supply recipe and production-operation concepts~\cite{isa88_2010,isa95_2025}. Run-to-run control motivates using observations in subsequent decisions, and human--automation research motivates separating proposal and approval~\cite{moyne2001r2r,parasuraman2000model}. The distinction here is the responsibility of the runtime: resolve an intent to the active scene and its direction-specific mappings, decide whether the operation may proceed, and return evidence without conflating command, readback, and process observation. The comparison is about where these responsibilities sit, not whether another system could implement them.
